%%% This is a tentative paragraph written by Pedro on August 11th to serve as an introduction only for Chapter 2. To be discussed with the rest of the group. 

The Publication Strategy Task Force was convened by the JUNO Executive Board in the summer 2024 with the primary goal of studying our experiment's sensitivity to the Neutrino Mass Ordering (NMO) when using external constraints on $\Delta m^2_{31}$, most notably from accelerator experiments. The task force's preliminary findings on this topic have been reported in several presentations, including a summary in Ref.~\cite{docdb:jantfreport}. Our results suggest that JUNO has a good chance of making a $\geq 3\sigma$ significant claim on the NMO with about 1 year of data when using existing external constraints from accelerator experiments, depending on the value of $\Delta m^2_{31}$ that JUNO measures. A dedicated technote presenting these results is in preparation. 

Informed by these findings, the collaboration plans to release a first result on the $\sin^2 \theta_{21}$, $\Delta m^2_{21}$ and possibly $\Delta m^2_{31}$ oscillation parameters with a low-statistics sample collected over roughly 50 days of data taking. This sample is expected to be large enough to achieve world-leading precision on the first two parameters and potentially on the last one as well. This result will be followed by another one on about roughly 1 year of data and incorporating external constraints, with the aim of making a significant claim on the NMO. Given the immediacy of the first result, the task force focused its attention during the last few months to studying the challenges and pitfalls involved in a low-statistics measurement of the oscillation parameters. The findings from this process are the focus of this technote and are being released so that the various analysis groups can take them into account.

The rest of this technote is organized as follows. \autoref{sec:2_bad_minima} discusses the complications arising from the large fluctuations expected in low-statistics data sets, namely the occurrence of bad minima and disjointed confidence intervals for $\Delta m^2_{31}$. \autoref{sec:coverage} addresses statistical coverage and shows that a Feldman–Cousins (or similar) approach is required to obtain proper confidence intervals, adding complexity to the analysis. \autoref{sec:solar_params} shows that these challenges are much milder for the solar parameters $\sin^2\theta_{12}$ and $\Delta m^2_{21}$.  \autoref{sec:dm231constraints} describes a possible mitigation, namely the incorporation of an external constraint on $\Delta m^2_{31}$. \autoref{sec:dybconstraints} presents the sensitivity to the oscillation parameters when using constraints from Daya Bay on the spectral shape and net flux of reactor antineutrinos, rather than from the TAO satellite detector. Finally, \autoref{sec:nmoimpact} studies the impact a low-statistics measurement of $\Delta m^2_{31}$ from JUNO can have on ongoing long-baseline neutrino experiments. 

